\documentclass[11pt]{article}

\usepackage[a4paper,margin=1in]{geometry}
\usepackage[T1]{fontenc}
\usepackage[utf8]{inputenc}
\usepackage[english]{babel}
\usepackage{amsmath,amssymb,amsthm}
\usepackage{booktabs}
\usepackage{hyperref}
\usepackage{microtype}

\newtheorem*{claim}{Claim}
\newtheorem*{remark}{Remark}

\title{From MLPs to transformers:\\
       where universality changes}
\author{}
\date{\today}

\begin{document}
\maketitle

\begin{abstract}
The transition from an MLP to a transformer is often described by saying
that transformers are more powerful or scale better. A more precise
architectural distinction is that the sequence length moves from the shape
of a flattened input weight matrix to the structure of the input itself.
Parameter sharing and attention make this possible. Universal approximation
theorems explain the resulting expressive capacity for each fixed sequence
length; they do not establish a single set of universal weights for all lengths.
\end{abstract}

\section{The starting point}

\paragraph{Cybenko/Hornik (1989, 1991).}
Under the relevant activation assumptions, every continuous
$f:K\subset\mathbb R^n\to\mathbb R^m$ on a compact set can be
approximated arbitrarily closely by an MLP with one hidden layer.
For hidden width $w$, the input weight matrix has $wn$ entries.

\paragraph{Yun, Bhojanapalli, Rawat, Reddi, Kumar (2020).}
For each fixed sequence length $T$, transformers with trainable positional
encodings can approximate any continuous
$f:K\subset\mathbb R^{d\times T}\to\mathbb R^{d\times T}$ on a compact
domain, in $L^p$ for $1\le p<\infty$. The numbers of heads, head size,
and feed-forward hidden width can be fixed at $(h,m,r)=(2,1,4)$,
independently of $T$ and $d$. The construction can increase its depth.
\footnote{The theorem uses a simplified transformer block without layer
normalization. These small values of $(h,m,r)$ do not mean that the
total parameter count is independent of $T$ or $d$.}

\section{The architectural shift}

\begin{claim}[the structural distinction]
For a flattened-input MLP, the input weight matrix
$W\in\mathbb R^{w\times(Td)}$ changes shape when $T$ changes.
For a transformer, the same attention and token-wise feed-forward
projection matrices apply at each position. Their shapes do not depend
on the sequence length $T$.
\end{claim}

In this sense, a dimension of the flattened input becomes a property
of the input sequence:
\begin{equation}
\underbrace{Td}_{\text{MLP input-layer shape}}
\quad\longrightarrow\quad
\underbrace{T}_{\text{sequence structure}}.
\end{equation}

This permits shared transformer blocks to process different sequence
lengths. It is an architectural statement, distinct from the quantifiers
in the approximation theorem:
\[
  \text{for every fixed }T,\ f,\ \epsilon>0,
  \quad\text{there exists an approximating network }g_{T,f,\epsilon}.
\]
The weights and depth may depend on $T$, $f$, and $\epsilon$.
Moreover, the trainable positional encoding in Yun is a
$d\times T$ matrix and therefore has a parameter count that depends on $T$.

\section{The mechanism: contextual mapping and the division of roles}

The constructive proof explains how the two types of layers divide
the work.

\begin{claim}[division of roles]
Self-attention layers compute contextual mappings. Token-wise
feed-forward layers map the resulting contextual identifiers to target values.
\end{claim}

A \emph{contextual mapping} associates each token with a scalar
identifier that depends on the full input sequence. On the subset of
discretized inputs used in the proof, these identifiers distinguish
positions within an input and tokens across different inputs, subject
to the relevant permutation symmetry. For example, occurrences of
``I'' in ``I am happy'' and ``I am Bob'' can receive different identifiers.

Yun constructs such a mapping on almost all points of a grid in
$[0,1]^{d\times T}$ using a sequence of self-attention layers.
Each layer uses only pairwise interactions and shares its projections
across positions, yet their composition encodes the full context.
Feed-forward layers then map the distinguishable identifiers to the
required outputs.

The contextual mapping explains how parameter sharing can coexist
with universal approximation. The architectural shift describes the
structure; the constructive proof supplies an expressive-capacity result.

\paragraph{Related empirical evidence.}
Yun also replaces attention layers with bilinear projections or
depthwise separable convolutions. Replacing the first one or two
attention layers in $\text{BERT}_{\text{BASE}}$ with convolutions can
improve SQuAD and MNLI results, while replacing more layers hurts.
These observations are consistent with cheap layers supplying some
contextual information. They do not by themselves prove which
mechanisms account for the remaining advantage of attention.

\section{The accompanying mechanisms}

\begin{enumerate}
\item \textbf{Parameter sharing across positions.} The same
$(W_Q,W_K,W_V)$ apply at each position. Their shapes remain fixed as
the number of positions changes.

\item \textbf{Permutation equivariance before positional encoding.}
Self-attention respects permutations of tokens. Positional encoding
provides information about order. Without it, the universal
approximation result is restricted to permutation-equivariant functions.

\item \textbf{Attention routing.} With full attention, each position
can access every other position in one layer. A causal mask changes
which interactions are available; the theorem discussed here concerns
the paper's full-attention architecture.
\end{enumerate}

\section{What the theorem does not establish}

\begin{remark}[computational cost]
Dense attention conventionally forms $O(T^2)$ pairwise scores.
Efficient implementations can reduce working memory, and sparse or
linear attention changes the computation. Universality itself gives
no guarantee of practical efficiency.
\end{remark}

\begin{remark}[the approximation resource: width and depth]
In the classical shallow-MLP construction, width supplies increasing
approximation capacity. Yun fixes $(h,m,r)=(2,1,4)$ and pays with depth.
For example, the contextual-mapping construction uses
$\delta^{-d}+1$ attention layers for grid spacing $\delta$.
Fixed local widths therefore do not imply a small approximating network.
\end{remark}

\begin{remark}[universal approximation and Turing completeness]
Yun concerns approximation of functions on a continuum. Results about
Turing completeness concern computation and depend on assumptions
about precision and available memory. They answer different questions.
\end{remark}

\begin{remark}[compact domains]
The positional-encoding result applies on a compact domain.
It does not guarantee uniform approximation over an unbounded input space.
\end{remark}

\begin{remark}[generalization to longer inputs]
The existence of an approximator for each fixed $T$ does not establish
that a trained network will generalize from $T\le T_{\text{train}}$
to much longer inputs. Length generalization is a separate question.
\end{remark}

\begin{remark}[positional information]
The paper's extension beyond permutation-equivariant targets uses
trainable positional encodings. The encoding is a substantive assumption
in that result.
\end{remark}

\section{Comparison by property}

\begin{center}
\small
\begin{tabular}{p{0.25\linewidth}p{0.30\linewidth}p{0.33\linewidth}}
\toprule
Property & Flattened-input MLP & Transformer \\
\midrule
Approximation domain & Fixed $\mathbb R^n\to\mathbb R^m$
  & Fixed $\mathbb R^{d\times T}\to\mathbb R^{d\times T}$ \\
Input projections & Depend on flattened input size
  & Shared across positions \\
Permutation equivariance & No built-in guarantee
  & Present before positional encoding \\
Approximation resource & Width in the shallow construction
  & Depth with fixed $(h,m,r)$ \\
Long-range access & Through dense input connections
  & Through full attention \\
Variable input length & Requires adapting the input interface
  & Supported by shared blocks, subject to positional design \\
Dense interaction cost & Depends on layer widths
  & $O(T^2)$ attention scores \\
\bottomrule
\end{tabular}
\end{center}

\section{The same idea in other architectures}

Sharing local parameters across a structural dimension of the input
is a broader architectural idea:

\begin{itemize}
\item \textbf{CNNs:} convolution shares parameters across spatial
positions and can operate on different spatial extents.
\item \textbf{Graph neural networks:} shared message-passing rules
operate on graphs with different numbers of vertices, with permutation
equivariance over vertex labels.
\item \textbf{Deep Sets:} shared element-wise maps and symmetric
aggregation support permutation-invariant functions of sets.
\end{itemize}

Each architecture requires its own expressive-capacity theorem.
Sharing parameters across a structural dimension does not automatically
establish universal approximation or generalization across input sizes.

\section{Conclusion}

The architectural distinction is precise: transformer projections are
shared across token positions, so their shapes need not grow with
sequence length. Yun explains how attention and feed-forward layers
can nevertheless approximate a rich class of functions for each fixed
length. The fixed values $(h,m,r)=(2,1,4)$ describe local widths;
the network's depth, weights, and positional encoding remain part of
the approximation construction.

\begin{thebibliography}{9}
\bibitem{cybenko1989} Cybenko, G.\ \emph{Approximation by superpositions
of a sigmoidal function}. Mathematics of Control, Signals and Systems,
1989.

\bibitem{hornik1991} Hornik, K.\ \emph{Approximation capabilities of
multilayer feedforward networks}. Neural Networks, 1991.

\bibitem{yun2020} Yun, C., Bhojanapalli, S., Rawat, A., Reddi, S.,
Kumar, S.\ \emph{Are Transformers Universal Approximators of
Sequence-to-Sequence Functions?}. ICLR 2020. arXiv:1912.10077.

\bibitem{zaheer2017} Zaheer, M.\ et al.\ \emph{Deep Sets}. NeurIPS 2017.

\bibitem{perez2019} P\'erez, J., Marinkovi\'c, J., Barcel\'o, P.\
\emph{On the Turing Completeness of Modern Neural Network Architectures}.
ICLR 2019.
\end{thebibliography}

\end{document}
